\documentclass[10pt,a4paper]{article}
\usepackage[english,spanish]{babel}
\usepackage{indentfirst}
\usepackage{anysize} % Soporte para el comando \marginsize
%\marginsize{1.5cm}{1.5cm}{0.5cm}{1cm}
\marginsize{2,5cm}{1,8cm}{4cm}{1,7cm}
\usepackage[psamsfonts]{amssymb}
\usepackage{amssymb}
\usepackage{amsfonts}
\usepackage{amsmath}
\usepackage{amsthm}
\usepackage{multicol}
\renewcommand{\thepage}{}
\columnsep=7mm

%%%%%%%%%%%%%%%%%%%%%%%%%%%%%%%%%%%%%%%%
\newtheorem{definicion}{Definici\'on}[section]
\newtheorem{teorema}{Teorema}[section]
\newtheorem{prueba}{Prueba}[section]
\newtheorem{prueba*}{Prueba}[section]
\newtheorem{corolario}{Corolario}[section]
\newtheorem{observacion}{Observaci\'on}[section]
\newtheorem{lema}{Lema}[section]
\newtheorem{ejemplo}{Ejemplo}[section]
\newtheorem{solucion*}{Soluci\'on}[section]
\newtheorem{algoritmo}{Algoritmo}[section]
\newtheorem{proposicion}{Proposici\'on}[section]

\linespread{1.4} \sloppy

\newcommand{\R}{\mathbf{R}}
\newcommand{\N}{\mathbf{N}}
\newcommand{\C}{\mathbb{C}}
\newcommand{\Lr}{\mathcal{L}}
\newcommand{\fc}{\displaystyle\frac}
\newcommand{\ds}{\displaystyle}

\DeclareMathOperator{\Dom}{Dom}

%%%%%%%%%%%%%%%%%%%%%%%%%%%%%%%%%%%%%%%%

\renewcommand{\thefootnote}{\fnsymbol{footnote}}

\begin{document}
\begin{center}
 {\Large \textbf{Análisis de seguridad frente a ataques criptográficos en el Cifrado de Hill}}
\end{center}
\begin{center}
 Dire Daniel Espinal Pecho $1$, Leonel Mathias Puente Rojas $2$, Luis Augusto Pareja Gavilan $3$, Preston Rodrigo Cordova Bedon $4$\vskip12pt
{\it Facultad de Ciencias $1$, Universidad Nacional de Ingeniería $1$, e-mail: dire.espinal.p@uni.pe \\
Facultad de Ciencias $2$, Universidad Nacional de Ingeniería $2$, e-mail: leonel.puente.r@uni.pe \\
Facultad de Ciencias $3$, Universidad Nacional de Ingeniería $3$, e-mail: luis.pareja.g@uni.pe \\
Facultad de Ciencias $4$, Universidad Nacional de Ingeniería $4$, e-mail: preston.cordova.b@uni.pe }
\end{center}
%\maketitle 
\begin{quotation}
{\small
%\begin{abstract}
\begin{center}
Resumen	
\end{center}
En este trabajo de investigación se busca examinar los aspectos fundamentales sobre la encriptación mediante el cifrado de Hill, analizaremos su relación directa con el álgebra lineal y la aritmética modular, pondremos a prueba su resistencia contra ataques criptográficos y detallaremos algunas aplicaciones de computación cuántica.
%\end{abstract}
Palabras Claves: álgebra lineal, aritmética modular, computación cuántica y ataques criptográficos. \\
}\\
{\small
\hspace*{0.5cm}
\begin{center}
Resumen en ingl\'es
\end{center}
This research paper seeks to examine the fundamental aspects of encryption using the Hill cipher, analyze its direct relationship with linear algebra and modular arithmetic, test its resistance against cryptographic attacks, and detail some quantum computing applications.

Keywords: linear algebra, modular arithmetic, quantum computing and cryptographic attacks. \\ 
}
\end{quotation}
\begin{multicols}{2}
\begin{center}
{\large \bf 1. INTRODUCCI\'ON}
\end{center}

Desde los albores de la civilización, la necesidad de transmitir información de manera segura y privada ha sido un motor fundamental para el desarrollo tecnológico. La criptografía, disciplina encargada de cifrar mensajes para que solo el destinatario legítimo pueda interpretarlos, ha evolucionado de métodos mecánicos simples a complejos algoritmos matemáticos. En este contexto, el Cifrado de Hill, propuesto por el matemático estadounidense Lester Hill en 1929, representa un hito histórico al ser el primer sistema criptográfico práctico que utiliza conceptos avanzados de álgebra lineal.

El presente trabajo complementa el análisis algebraico del cifrado de Hill con su implementación y validación práctica. Se desarrolla un programa en C++ que ejecuta el cifrado y descifrado mediante operaciones modulares sobre matrices invertibles, comprobando cómo la existencia del inverso modular de la matriz clave condiciona la correcta recuperación del mensaje. Asimismo, se simula el sistema en IBM Composer para contrastar los resultados teóricos y prácticos, verificando la coherencia entre el modelo algebraico y su comportamiento real. Con esta doble aproximación se evalúa la seguridad del cifrado frente a ataques como la fuerza bruta y el known-plaintext attack, demostrando que su solidez depende de las propiedades estructurales del álgebra lineal en cuerpos finitos.

Un aspecto fundamental que distingue al Cifrado de Hill de otros métodos clásicos es su resistencia frente al análisis de frecuencias. Mientras que los cifrados de sustitución simple son vulnerables porque mantienen la frecuencia estadística de las letras del idioma, el método de Hill, al ser un cifrado polialfabético que opera sobre bloques (digramas, trigramas, etc.), logra ”difundir” la estructura del lenguaje a través de la matriz de transformación. Esto significa que una misma letra en el texto original puede convertirse en diferentes caracteres en el criptograma, dependiendo de su posición en el bloque y de los valores de la matriz clave, lo que obliga al criptoanalista a enfrentarse a un sistema de ecuaciones lineales mucho más complejo para vulnerar la seguridad del mensaje.

Asimismo, el Cifrado de Hill ayuda a entender cómo funcionan algunos sistemas de cifrado actuales, donde la información se trabaja en bloques de datos en lugar de analizar cada letra por separado. Este tipo de enfoque permite observar de manera más clara la relación entre las matemáticas y la seguridad de la información. A partir de ello, se puede reconocer la importancia de herramientas como las matrices en el procesamiento de datos. Además, el estudio de este método contribuye a fortalecer la comprensión de conceptos básicos del álgebra lineal, como los determinantes y las matrices inversas. De esta manera, se busca resaltar la utilidad de estos conocimientos en situaciones aplicadas, más allá de su uso teórico en el aula.

\begin{center}
{\large \bf 2. CONCEPTOS PREVIOS}
\end{center}

\subsection*{2.1. El Espacio Vectorial $\mathbb{Z}_m$}

El texto plano se segmenta en bloques de longitud $m$, donde cada bloque se representa como un vector columna $P \in \mathbb{Z}_m$. La elección de la dimensión $m$ define el tamaño del bloque y, por extensión, la complejidad del sistema. El espacio de trabajo es el producto cartesiano de $m$ copias del anillo $\mathbb{Z}_n$ (donde $m$ es el orden del alfabeto). Toda operación se realiza bajo la métrica de congruencia módulo $n$.

\subsection*{2.2. El Concepto de Anillo $\mathbb{Z}_n$}

En matemáticas formales, el cifrado de Hill opera en el anillo de enteros módulo $n$, donde $n$ es el tamaño del alfabeto empleado. Decimos que dos números $a$ y $b$ son congruentes módulo $m$ si su diferencia es un múltiplo exacto de $n$:
\[
a \equiv b \pmod{n} \Leftrightarrow \exists k \in \mathbb{Z} \text{ tal que } a - b = kn.
\]
Esto garantiza que cualquier resultado de una multiplicación de matrices se mantenga dentro del rango del alfabeto $[0, n-1]$, ya que si lo sobrepasa, el resultado se reduce módulo $n$ para identificar la letra correspondiente.

\subsection*{2.3. La Matriz de Transformación en $M_{m\times m}(\mathbb{Z})$}

La clave del sistema es una matriz cuadrada $K$ de orden $m\times m$, cuyas entradas $k_{ij}$ pertenecen al anillo $\mathbb{Z}_m$. La operación de cifrado se define como la multiplicación matriz-vector $K \cdot P$, resultando en un vector de criptograma $C$ de la misma dimensión $m$:
\[
C \equiv K \times P \pmod{n}, \quad \text{donde } C, P \in \mathbb{Z}_n.
\]
Esto es válido porque $K^{-1} \in M_{m\times m}(\mathbb{Z}_n)$ y $P \in \mathbb{Z}_n^m$, por lo que cada entrada del vector resultante es una combinación lineal de enteros reducida módulo $n$, garantizando que $C \in \mathbb{Z}_n^m$. Es decir, la transformación lineal $T(P)=K\cdot P$ es un endomorfismo del módulo $\mathbb{Z}_n^m$ sobre sí mismo. Para que el proceso sea reversible, la matriz $K$ debe ser un elemento del Grupo General Lineal de grado $m$ sobre $\mathbb{Z}_n$, denotado como $GL(m,\mathbb{Z}_n)$. Esto garantiza que la transformación lineal sea biyectiva.

\subsection*{2.4. Inverso Multiplicativo}

En la aritmética modular y en especial en el Anillo $\mathbb{Z}_m$, se define al inverso multiplicativo de un número $a$ como aquel $a^{-1}$ tal que $a \times a^{-1} \equiv 1 \pmod{m}$. Tal inverso existe solo si $\operatorname{mcd}(a,m)=1$ (son coprimos).

\subsection*{2.5. Condiciones de Invertibilidad y Dimensión}

La existencia de la matriz de descifrado $K^{-1} \in M_{m\times m}(\mathbb{Z}_n)$ depende de una condición necesaria para que sea invertible en dicho módulo $n$:
\[
\operatorname{mcd}(\det(K), n) = 1.
\]
Por el Teorema de Bézout, esta condición garantiza la existencia de $\det(K)^{-1}$ en $\mathbb{Z}_n$, permitiendo construir:
\[
K^{-1} = \frac{1}{\det(K)} \times \operatorname{adj}(K).
\]
En este caso no se usa $\det(K) \neq 0$, ya que es redundante. Si $\operatorname{mcd}(\det(K), n)=1$ entonces $\det(K)$ no puede ser $0$ en $\mathbb{Z}_n$, pues $\operatorname{mcd}(0,n)=n \neq 1$ para todo $n>1$. La condición $\det(K) \neq 0$ queda así contenida en la condición de coprimalidad.

Ejemplo: En $\mathbb{Z}_{26}$, si $\det(K)=2$, entonces $\operatorname{mcd}(2,26)=2 \neq 1$, por lo que no existe $K^{-1}$ en módulo 26.

La existencia de $K^{-1}$ garantiza que $T(P)=K\times P$ es un automorfismo de $\mathbb{Z}_n^m$, es decir, una transformación biyectiva. En consecuencia, el mensaje original se recupera de forma única mediante:
\[
P = K^{-1} \times C \pmod{n}.
\]

\subsection*{2.6. Independencia Lineal y Criptoanálisis}

En un ataque de texto plano conocido, el objetivo es determinar las $m^2$ incógnitas de la matriz $K$. Para ello, se requiere capturar al menos $m$ vectores linealmente independientes en $\mathbb{Z}_n^m$. Si los vectores expuestos $\{P_1, P_2, \dots, P_m\}$ no forman una base para el espacio $\mathbb{Z}_n^m$ (es decir, el rango de la matriz formada por ellos es menor a $m$), el sistema de ecuaciones lineales será indeterminado, impidiendo la recuperación única de la clave.

Una vez que los vectores expuestos formen una base, se agrupan en una matriz cuadrada $P$ y los vectores de texto cifrado correspondientes en otra matriz $C$. Las letras de ambas se organizan como vectores columna:
\[
K \times P = C \pmod{n}.
\]
Se deben cumplir las condiciones de invertibilidad para la matriz $P$: $\operatorname{mcd}(\det(P), n)=1$ y $\det(P) \neq 0 \pmod{n}$. Calculamos la inversa modular de $P$ mediante:
\[
P^{-1} = (\det(P))^{-1} \cdot \operatorname{adj}(P) \pmod{n}.
\]
Al multiplicar ambos lados de la ecuación por $P^{-1}$, se despeja la clave:
\[
K \equiv C \times P^{-1} \pmod{n}.
\]
Así se puede descifrar cualquier mensaje, rompiendo el código.

\subsection*{2.7. La relación directa con la matriz llave}

Para descifrar el mensaje se necesita la inversa de la matriz clave. En álgebra lineal, la fórmula es:
\[
K^{-1} = \frac{1}{\det(K)} \times \operatorname{adj}(K).
\]
En el cifrado de Hill, esa fracción se convierte en el inverso multiplicativo modular del determinante. Los pasos matemáticos para la validez de una clave son:
\begin{enumerate}
\item Calcular el determinante $d = \det(K)$.
\item Verificar si son coprimos: calcular $\operatorname{mcd}(d, n)$. Si no es 1, la matriz no sirve como clave.
\item Hallar el inverso modular de $d$ y construir $K^{-1}$.
\end{enumerate}

Ejemplo: Sea $K$ una matriz $3\times 3$:
\[
K = \begin{bmatrix} 1 & 2 & 3 \\ 0 & 1 & 4 \\ 5 & 6 & 0 \end{bmatrix}.
\]
Se tiene $\det(K)=1$, $\operatorname{mcd}(1,26)=1$, y como $\det(K)^{-1}=1$ en $\mathbb{Z}_{26}$,
\[
K^{-1} = \operatorname{adj}(K) \pmod{26} = \begin{bmatrix} 2 & 18 & 5 \\ 20 & 11 & 22 \\ 21 & 4 & 1 \end{bmatrix}.
\]
Al pasar de una matriz $2\times2$ a una $3\times3$, el espacio de claves crece de $26^4$ a $26^9$, incrementando drásticamente. Este crecimiento será importante para un ataque de fuerza bruta.

\subsection*{2.8. Computación clásica y cuántica aplicada al criptoanálisis del Cifrado de Hill}

\subsubsection*{2.8.1. Computación Cuántica}

Un \textbf{qubit} es la unidad básica de información cuántica. A diferencia del bit clásico, puede estar en superposición $\alpha|0\rangle + \beta|1\rangle$, lo que permite a un computador cuántico evaluar múltiples claves del Cifrado de Hill simultáneamente:
\[
|\psi\rangle = e^{i\gamma} \left( \cos \frac{\theta}{2} |0\rangle + e^{i\varphi} \sin \frac{\theta}{2} |1\rangle \right).
\]
La superposición permite representar todos los estados del espacio de claves $26^{m^2}$ al mismo tiempo, condición que explota el algoritmo de Grover. El entrelazamiento cuántico es un fenómeno en el que dos qubits se entrelazan de tal forma que el estado de una partícula no puede describirse independientemente del estado de la otra, sin importar la distancia que las separe.

\subsubsection*{2.8.2. Algoritmos cuánticos aplicados a la criptografía}

El \textbf{algoritmo de Grover} permite resolver el problema de búsqueda por fuerza bruta con complejidad $O(\sqrt{N})$ en lugar de $O(N)$ consultas. Utiliza superposición, inversión de fase (oráculo) y difusión para amplificar la probabilidad del resultado correcto. Dado que el espacio de claves de Hill es $N = 26^{m^2}$, Grover reduce la búsqueda a $O(\sqrt{N}) = O(26^{m^2/2})$ consultas. Para $m=2$, esto pasa de $456\,976$ a aproximadamente $676$ operaciones cuánticas.

\begin{center}
{\large \bf 3. AN\'ALISIS}
\end{center}

\subsection*{3.1. Modelo Matemático}

\subsubsection*{3.1.1. Desarrollo del cifrado}

Paso 1) Escoger una matriz llave invertible módulo 26. Sea $K$ de orden 2:
\[
K = \begin{pmatrix} 3 & 3 \\ 2 & 5 \end{pmatrix}.
\]
Calculamos $\det(K)=9$ y verificamos $\operatorname{mcd}(9,26)=1$; por Euclides se comprueba que 9 y 26 son coprimos.

Paso 2) Escoger un mensaje para cifrar y convertirlo en vectores por bloques. Sea el mensaje \texttt{LINEAL}. Asignamos valores: $L=11, I=8, N=13, E=4, A=0, L=11$. Formamos bloques de dos:
\[
P_1 = \begin{pmatrix} 11 \\ 8 \end{pmatrix}, \quad
P_2 = \begin{pmatrix} 13 \\ 4 \end{pmatrix}, \quad
P_3 = \begin{pmatrix} 0 \\ 11 \end{pmatrix}.
\]

Paso 3) Multiplicar cada bloque por la llave:
\[
K P_1 = \begin{pmatrix} 57 \\ 62 \end{pmatrix},\quad
K P_2 = \begin{pmatrix} 51 \\ 46 \end{pmatrix},\quad
K P_3 = \begin{pmatrix} 33 \\ 55 \end{pmatrix}.
\]

Paso 4) Reducir módulo 26:
\[
C_1 = \begin{pmatrix} 5 \\ 10 \end{pmatrix},\quad
C_2 = \begin{pmatrix} 25 \\ 20 \end{pmatrix},\quad
C_3 = \begin{pmatrix} 7 \\ 3 \end{pmatrix}.
\]

Paso 5) Convertir a letras: $C_1 = FK$, $C_2 = ZU$, $C_3 = HD$. El mensaje cifrado es \texttt{FKZUH D} (notar que la última D proviene del tercer bloque).

\subsubsection*{3.1.2. Desarrollo del descifrado}

\paragraph{I) Algebraico:}
Paso 1) Convertir el mensaje cifrado \texttt{FKZUH D} a vectores: $F=5, K=10, Z=25, U=20, H=7, D=3$. Bloques:
\[
C_1 = \begin{pmatrix} 5 \\ 10 \end{pmatrix},\quad
C_2 = \begin{pmatrix} 25 \\ 20 \end{pmatrix},\quad
C_3 = \begin{pmatrix} 7 \\ 3 \end{pmatrix}.
\]

Paso 2) Hallar la inversa de $K$:
\[
\operatorname{adj}(K) = \begin{pmatrix} 5 & -3 \\ -2 & 3 \end{pmatrix}, \quad \det(K)=9, \quad 9^{-1} \pmod{26} = 3.
\]
Entonces:
\[
K^{-1} = 3 \cdot \begin{pmatrix} 5 & -3 \\ -2 & 3 \end{pmatrix} \pmod{26} = \begin{pmatrix} 15 & 17 \\ 20 & 9 \end{pmatrix}.
\]

Paso 3) Multiplicar cada bloque cifrado por $K^{-1}$:
\[
P_1 = \begin{pmatrix} 245 \\ 190 \end{pmatrix}, \quad
P_2 = \begin{pmatrix} 715 \\ 680 \end{pmatrix}, \quad
P_3 = \begin{pmatrix} 156 \\ 167 \end{pmatrix}.
\]

Paso 4) Reducir módulo 26:
\[
P_1 = \begin{pmatrix} 11 \\ 8 \end{pmatrix},\quad
P_2 = \begin{pmatrix} 13 \\ 4 \end{pmatrix},\quad
P_3 = \begin{pmatrix} 0 \\ 11 \end{pmatrix}.
\]

Paso 5) Convertir a letras: $LI, NE, AL$, recuperando \texttt{LINEAL}.

\paragraph{II) Known-Plaintext Attack}

Supongamos que el atacante conoce el texto plano y el cifrado correspondiente. Se busca recuperar $K$ a partir de la relación $C = K \times P \pmod{26}$. Si se tienen suficientes bloques linealmente independientes, se forma la matriz $P$ y se calcula $K = C \times P^{-1} \pmod{26}$.

En el ejemplo, se tienen $P_1, P_2, P_3$. Se prueba con $P_1$ y $P_3$ porque $P_1$ y $P_2$ resultan linealmente dependientes (determinante 18, no coprimo con 26). Con $P_1$ y $P_3$:
\[
P = \begin{pmatrix} 11 & 0 \\ 8 & 11 \end{pmatrix}, \quad \det(P)=121 \equiv 17 \pmod{26}, \quad \operatorname{mcd}(17,26)=1.
\]
Se calcula $P^{-1}$ con $\det(P)^{-1}=23$:
\[
P^{-1} = 23 \cdot \begin{pmatrix} 11 & 0 \\ 18 & 11 \end{pmatrix} \pmod{26} = \begin{pmatrix} 19 & 0 \\ 24 & 19 \end{pmatrix}.
\]
La matriz $C$ correspondiente es:
\[
C = \begin{pmatrix} 5 & 7 \\ 10 & 3 \end{pmatrix}.
\]
Entonces:
\[
K = C \times P^{-1} \pmod{26} = \begin{pmatrix} 3 & 3 \\ 2 & 5 \end{pmatrix},
\]
recuperando la clave.

\paragraph{III) Brute Force Attack}

El atacante solo tiene el texto cifrado y asume un cifrado de Hill con bloques de 2. El espacio de claves es $26^4 = 456\,976$ matrices. Se escribe un script que genera cada matriz, verifica la invertibilidad ($\operatorname{mcd}(\det(K),26)=1$), calcula su inversa, descifra y comprueba si el texto resultante tiene sentido (por ejemplo, contiene un fragmento esperado). Este ataque es tedioso para una persona, pero factible con computadora.

\paragraph{IV) IBM Composer}

Se implementa el algoritmo de Grover en IBM Quantum Composer para buscar una clave de 3 bits (espacio de 8 claves). Se inicializan 4 qubits (3 de búsqueda y 1 ancilla), se aplican compuertas Hadamard para crear superposición, se implementa el oráculo que marca el estado objetivo (por ejemplo, $|101\rangle$) y luego el operador de difusión. Tras una iteración, la probabilidad del estado marcado aumenta de 12.5\% a aproximadamente 78\%, demostrando la ventaja cuántica.

\subsection*{3.2. Análisis}

A partir de las simulaciones y cálculos, se desprenden los siguientes puntos:

\begin{itemize}
\item \textbf{Restricciones matemáticas del espacio de claves ($\mathbb{Z}_{26}$):} La condición de invertibilidad es $\operatorname{mcd}(\det(K),26)=1$. Si no se cumple, la clave no sirve.

\item \textbf{Linealidad del cifrado:} El cifrado es lineal, lo que lo hace vulnerable a ataques de texto plano conocido si se obtienen bloques linealmente independientes.

\item \textbf{Condiciones de éxito en el ataque:} Es obligatorio que los bloques interceptados formen una base válida; de lo contrario, el sistema es indeterminado.

\item \textbf{Impacto del algoritmo de Grover:} La complejidad clásica $O(N)$ se reduce a $O(\sqrt{N})$ con computación cuántica. Para $m=2$, de $456\,976$ a $676$ intentos.
\end{itemize}

\begin{center}
{\large \bf 4. OBSERVACIONES}
\end{center}

Se observa que la seguridad del Cifrado de Hill depende críticamente del tamaño de la matriz y de la elección del alfabeto. Para $m=2$, el espacio de claves es manejable por fuerza bruta clásica, pero para $m$ grande (por ejemplo, $m=10$) el espacio $26^{100}$ es inabordable. Sin embargo, la linealidad sigue siendo una debilidad fundamental. La implementación en C++ y en IBM Composer confirma los resultados teóricos y muestra la viabilidad práctica de los ataques.

\begin{center}
{\large \bf 5. CONCLUSIONES}
\end{center}

\begin{itemize}
\item La seguridad del Cifrado de Hill no se basa únicamente en mantener la matriz clave $K$ en secreto, sino en verificar que pertenezca al grupo lineal general y en el tamaño de esta matriz.

\item El éxito del ataque de texto plano conocido depende principalmente de la independencia lineal de los mensajes interceptados; de no serlo, el sistema genera infinitas soluciones.

\item A través del Algoritmo de Grover, se reduce drásticamente el tiempo y costo computacional de un ataque de fuerza bruta clásica de orden $O(N)$ a $O(\sqrt{N})$, forzando a usar matrices de mayor dimensión para contrarrestarlo.
\end{itemize}

%\begin{center}
%{\large \bf Agradecimientos}
%\end{center}
%Los autores agradecen a las autoridades de la Facultad de Ciencias de la Universidad Nacional de 
%Ingenier\'{\i}a por su apoyo.
%%\begin{center}
%%{\large \bf Apendice: }
%%\end{center}

\end{multicols}

\begin{center}
 -----------------------------------------------------------------------------
\end{center}
\begin{multicols}{2}
\begin{list}{}{\setlength{\topsep}{0mm}\setlength{\itemsep}{0mm}%
\setlength{\parsep}{0mm}\setlength{\leftmargin}{4mm}}
%
%------------------------------------- References --------------------
\small
\item[1.] Hill, L. S. (1929). Cryptography in an Algebraic Alphabet. The American Mathematical Monthly, 36(5), 306-312. \url{https://doi.org/10.2307/2299749}
\item[2.] Anton, H., y Rorres, C. (2014). Elementary Linear Algebra: Applications Version (11th ed.). Wiley.
\item[3.] Dooley, John (January 1, 2018). "10.1 The Shoulders of Giants: Friedman, Hill, and Shannon". History of Cryptography and Cryptanalysis: Codes, Ciphers, and Their Algorithms. Springer. p. 167. ISBN 978-3-319-90442-9.
\item[4.] Lester S. Hill, Concerning Certain Linear Transformation Apparatus of Cryptography, The American Mathematical Monthly Vol.38, 1931, pp. 135-154.
\item[5.] Stallings, W. (2017). Cryptography and network security: Principles and practice (7th ed.). Pearson, pp. 86-107.
\item[6.] Singh, S. (1999). The Code Book: The science of secrecy from ancient Egypt to quantum cryptography. Anchor Books.
\item[7.] Ibanez, R. (11 de enero de 2017). Cultura científica. \url{https://culturacientifica.com/2017/01/11/critopatrices-cifrado-hill/}
\item[8.] Jacques-Garcia, F. A., Uribe-Mejía, D., Macías-Bobadilla, G., y Chaparro-Sánchez, R. (2022). On modular inverse matrices a computation.
\item[9.] \url{https://www.ibm.com/es-es/think/topics/qubit}
\item[10.] Thomas G. Wong (2022). Introduction to classical and quantum computing. pp 299-303.
\item[11.] Michael A. Nielsen, M. A., y Isaac L. Chuang, I. L. (2010). Quantum computation and quantum information (10th anniversary ed.) pp 13-20.
\item[12.] Sipser, M. (2012). Introduction to the theory of computation (3rd ed.). Cengage Learning.
%---------------------------------------------------------------------
%
\end{list}
\end{multicols}
\end{document}
